
\section{Planteamiento del problema}

\subsection{Planteamiento del problema}

Los indicadores de continuidad del servicio SAIDI (Índice de Duración Promedio de Interrupción del Sistema) y SAIFI (Índice de Frecuencia Promedio de Interrupción del Sistema) en la empresa de distribución objeto de estudio presentan dificultades estructurales para su mejora sostenida. Si bien existen técnicas estadísticas y de aprendizaje automático ampliamente documentadas para el análisis de interrupciones en sistemas eléctricos, la principal limitación no radica en la ausencia de metodologías analíticas, sino en la calidad, estructura e interpretación de los datos disponibles, lo que dificulta el entrenamiento de modelos predictivos que sean simultáneamente confiables y explicativos \cite{ref1,ref4,ref3}. En particular, resulta complejo detectar, cuantificar y explicar no solo las causas de las interrupciones, sino también su duración y los tiempos de restauración del servicio.

Los principales obstáculos identificados son los siguientes:

\begin{itemize}
    \item \textbf{Fragmentación de las fuentes de información}: La operación del sistema requiere el uso de múltiples registros provenientes de plataformas heterogéneas como SCADA, OMS, GIS, CIS, bases de mantenimiento y fuentes meteorológicas externas. Estos datos se encuentran en distintos formatos, con estructuras disímiles y, en algunos casos, con ausencia de metadatos completos, lo que dificulta su integración sin pérdida de trazabilidad \cite{ref5,ref6}.
    \item \textbf{Problemas de sincronización temporal}: Los sistemas operativos no presentan una alineación temporal uniforme. Mientras algunos registran información en intervalos horarios, otros lo hacen en escalas de minutos o presentan series discontinuas. Esta falta de sincronización genera inconsistencias que afectan la estimación precisa de la duración de las interrupciones y de los tiempos de reposición del servicio \cite{ref5,ref7}.
    \item \textbf{Eventos extremos y ruido en los datos}: Los registros contienen alta variabilidad, errores de captura, información incompleta y eventos atípicos, como fenómenos climáticos extremos o fallas excepcionales, que suelen ser excluidos del análisis. Esta situación dificulta la identificación de patrones robustos tanto a nivel agregado como local \cite{ref1,ref8}.
    \item \textbf{Pérdida de detalle por agregación}: El uso de indicadores promediados en escalas anuales o geográficas amplias oculta el comportamiento específico de circuitos o zonas particulares, lo que limita la identificación de causas operativamente accionables \cite{ref4,ref7}.
    \item \textbf{Ausencia de estándares de datos}: La falta de catálogos de datos, definiciones homogéneas y métricas de integridad impide evaluar de manera sistemática si las fuentes disponibles son adecuadas para la construcción de modelos predictivos confiables \cite{ref5,ref6}.
    \item \textbf{Restricciones regulatorias}: Los organismos de control exigen que los resultados sean trazables y explicables. En consecuencia, no es suficiente emplear modelos con alto desempeño predictivo; estos deben ser interpretables y compatibles con los requerimientos de reporte regulatorio, lo que condiciona el tipo de metodologías que pueden adoptarse \cite{ref3,ref4}.
\end{itemize}

Estos obstáculos reducen la capacidad de la organización para explicar la variación de los indicadores SAIDI y SAIFI, priorizar actividades operativas y cumplir con los requerimientos regulatorios mediante herramientas predictivas \cite{ref3,ref8,ref5}. En términos cuantitativos, el SAIDI promedio de la empresa es de 22.43 horas por año, valor que se ubica un 4.64\% por debajo de la meta establecida por la CREG (23.522 horas por año). Para el indicador SAIFI, el promedio anual es de 6.5 eventos, significativamente inferior al objetivo regulatorio de 9 eventos por año, lo que representa una reducción del 27.78\% \cite{ref43}. Estos resultados sugieren que la empresa se encuentra dentro de los límites regulatorios; sin embargo, la gestión de dichos indicadores enfrenta importantes desafíos operativos e informativos.

Las compensaciones anuales asociadas a calidad del servicio ascienden aproximadamente a \$5 mil millones de pesos, y el análisis de datos históricos de cinco años evidencia que cerca del 2\% de los registros presenta información faltante o inconsistente, como tiempos de evento no registrados, campos incompletos o desajustes temporales entre sistemas. Si bien la Circular 063 de 2019 de la SSPD establece criterios técnicos para la exclusión válida de eventos asociados a factores externos, como fuerza mayor o fallas del Sistema Nacional de Transmisión, la presencia de problemas de calidad de datos limita la capacidad de construir modelos predictivos sólidos, diagnosticar con precisión las causas de las interrupciones y priorizar acciones operativas basadas en evidencia \cite{ref1,ref4,ref3,ref44}.

Este problema no es exclusivo de la empresa analizada. En el sector eléctrico colombiano, la calidad del servicio constituye un desafío estructural. De acuerdo con el Informe de Calidad del Servicio de Energía Eléctrica de la SSPD 2023, el SAIDI promedio nacional fue de 16.83 horas por usuario, con valores que oscilaron entre 4.5 horas y más de 50 horas dependiendo del operador y de la región geográfica \cite{ref9}. Esta dispersión refleja diferencias significativas en infraestructura, gestión operativa y exposición a condiciones climáticas extremas. Durante 2023, el 68\% de las interrupciones estuvo asociado a eventos externos, como contacto con vegetación, animales o terceros, mientras que el 24\% se atribuyó a fallas técnicas en equipos de red, evidenciando la naturaleza multicausal del problema \cite{ref4}.

El impacto regulatorio es considerable: en 2023, las compensaciones totales por incumplimiento de metas de calidad alcanzaron aproximadamente \$487\,000 millones a nivel nacional, siendo los indicadores de continuidad responsables del 73\% de este valor \cite{ref4}. Para operadores de red ubicados en zonas con alta exposición climática, como Norte de Santander, la problemática se intensifica. Según datos de la SSPD consolidados por la UPME, los eventos meteorológicos extremos incrementaron en un 35\% las interrupciones no programadas respecto al periodo 2019--2021 \cite{ref4}. En este contexto, la dispersión de información en sistemas heterogéneos dificulta la identificación de patrones causales, la priorización del mantenimiento preventivo y la detección temprana de circuitos críticos, limitando la capacidad de las empresas distribuidoras para anticipar deterioros en la continuidad del servicio y cumplir de manera consistente con los objetivos regulatorios establecidos por la CREG.
